\documentclass{article}
\usepackage[T1]{fontenc}
\usepackage{lmodern}
\usepackage{graphicx}
\usepackage{amsmath,amssymb}
\usepackage[a4paper,margin=2cm]{geometry}
\usepackage[absolute,overlay]{textpos}
\setlength{\TPHorizModule}{1mm}
\setlength{\TPVertModule}{1mm}
\usepackage[indonesian]{babel}
\usepackage{hyperref}
\setlength{\emergencystretch}{2em}
% unit: Halaman 5

\begin{document}

% === Halaman 5 (python/native) ===

5

Apa Steganografi itu?

• Dari Bahasa Yunani: steganos + graphien


“steganos”  (στεγανός):  tersembunyi


“graphien”  (γραφία)    : tulisan

 steganografi:  tulisan tersembunyi (covered writing)

• Steganography: ilmu dan seni menyembunyikan pesan rahasia dengan 

suatu cara sedemikian sehingga tidak seorang pun yang mencurigai 

keberadaan pesan tersebut.

   Tujuan steganografi: pesan tidak terdeteksi  keberadaannya

\end{document}
